% Exercises for Chapter 14 (American options by simulation). Statements only: solutions are kept privately.
% Notation as in the notes: exercise dates 0 = t_0 < ... < t_m = T, discounted payoffs Z_k,
% Snell envelope U, continuation values C_k, stopping times T_k, risk-neutral expectations E.

\begin{exercise}[Stopping a sequence of uniforms]\label{ex:am-secretary}
\easy Let $Z_1,\dots,Z_m$ be independent and uniform on $[0,1]$, $Z_0=0$,
$\F_k=\sigma(Z_1,\dots,Z_k)$, and no discounting. You may use \noteref{thm:am-dp}.
\begin{enumerate}[label=(\alph*)]
  \item Show that every continuation value is a constant, $C_k=v_{k+1}$, where $v_m=\frac12$ and
    $v_k=\frac12(1+v_{k+1}^2)$ for $1\le k<m$, and that $U_0=v_1$. Describe the optimal rule
    and compute $U_0$ for $m=2$ and $m=5$.
  \item Compute the perfect-foresight bound $\E[\max_kZ_k]$ and its gap for $m=2$ and $m=5$.
  \item Write down the martingale $M^*$ of \noteref{thm:am-dp}(v) and check directly, without
    \noteref{thm:am-duality}, that $\max_k(Z_k-M^*_k)=U_0$ for every outcome.
  \item Let $e_k=1-v_k$. Show that $e_k=e_{k+1}-\frac12e_{k+1}^2$ and deduce
    $1-U_0\le2/(m+3)$. Show that the perfect-foresight gap is at most $1/m$.
\end{enumerate}
\end{exercise}

\begin{exercise}[Bounds, and exercise at $t_0$]\label{ex:am-exercise-now}
\easy Consider a Bermudan put on the Black--Scholes stock, $Z_k=e^{-rt_k}(K-S_{t_k})^+$.
\begin{enumerate}[label=(\alph*)]
  \item Let $r\ge0$ and let $U_0(s)$ be the Bermudan price as a function of $S_0=s$, the
    dates being fixed. Show that $\max\big(K-s,P(0,s)\big)\le U_0(s)\le K$, where $P$ is the
    European put price, and that $U_0$ is convex and nonincreasing in $s$.
  \item Let $r>0$. Show that $C_0\le e^{-rt_1}K$, so that immediate exercise is optimal
    whenever $S_0\le K(1-e^{-rt_1})$. Evaluate this bound for the put of
    \noteref{tab:am-ls} with $\sigma=20\%$ and $T=1$, and compare it with the critical price
    at $t_0$, which you may compute with \texttt{bermudan\_grid}.
  \item The in-sample LSM estimate is $\max(Z_0,\hat c_0)$, where $\hat c_0$ is a mean of $N$
    cash flows. Suppose that the cash flows follow the optimal rule from $t_1$ on, so that
    $\E[\hat c_0]=C_0$. Show that $\E[\max(Z_0,\hat c_0)]\ge U_0$. If $Z_0=C_0$ and $\hat c_0$
    is normal with standard deviation $s$, compute the bias.
  \item How does the out-of-sample estimator of the notes treat $t_0$, and why is it not
    biased upwards there?
\end{enumerate}
\end{exercise}

\begin{exercise}[Two exercise dates in closed form]\label{ex:am-two-dates}
\hard Let a put be exercisable at $t_0=0$, $t_1=T/2$ and $t_2=T$, with $r>0$, $y\ge0$, and
let $P(\theta,s)$ be the Black--Scholes European put price with time to maturity $\theta$.
Write $N_2(a,b;\rho)$ for the distribution function of a standard bivariate normal vector
with correlation $\rho$.
\begin{enumerate}[label=(\alph*)]
  \item Show that $\varphi(s)=(K-s)-P(T/2,s)$ is strictly decreasing on $(0,\infty)$, with
    $\varphi(0+)>0>\varphi(K)$. Deduce that there is a unique $s^*\in(0,K)$ such that
    exercise at $t_1$ is optimal if and only if $S_{t_1}\le s^*$.
  \item Show that
    \[
      U_0=\max\Big(Z_0,\ e^{-rT/2}\E\big[(K-S_{T/2})\1_{\{S_{T/2}\le s^*\}}\big]
      +e^{-rT}\E\big[(K-S_T)^+\1_{\{S_{T/2}>s^*\}}\big]\Big),
    \]
    and write the two expectations with $N$ and $N_2(\cdot,\cdot;\rho)$, $\rho=\sqrt{1/2}$.
  \item Compute $s^*$ and $U_0$ for $S_0=K=40$, $r=6\%$, $y=0$, $\sigma=20\%$, $T=1$, and
    compare with \texttt{bermudan\_grid} with $m=2$ and with an out-of-sample LSM estimate.
\end{enumerate}
\end{exercise}

\begin{exercise}[A better control]\label{ex:am-stopped-control}
\medium Let $P(t,s)$ be the Black--Scholes price at $t$ of the European put with maturity
$T$ ($P(T,s)=(K-s)^+$) and $X_k=e^{-rt_k}P(t_k,S_{t_k})$.
\begin{enumerate}[label=(\alph*)]
  \item Show that $\E[X_\tau]=P(0,S_0)$ for every $\tau\in\mathcal T_0$.
  \item In the out-of-sample LSM estimator, use $X_{\hat\tau}$ as a control variate, with a
    coefficient estimated on the training paths. Show that the controlled estimator is
    still conditionally unbiased for the value of the policy.
  \item Implement it and compare its variance with that of the terminal control
    $e^{-rT}(K-S_T)^+$ of \noteref{exa:am-test-case}, on the put with $S_0=36$, $\sigma=20\%$,
    $T=1$ of \noteref{tab:am-ls} and on the test case. Explain the difference.
\end{enumerate}
\end{exercise}

\begin{exercise}[The European price as a basis function]\label{ex:am-european-basis}
\medium On the test case, \noteref{exa:am-test-case} found an LSM policy that exercises a third
of the paths although early exercise is worth almost nothing.
\begin{enumerate}[label=(\alph*)]
  \item What is the continuation value $c_k$ when early exercise is never optimal after
    $t_k$? Deduce a natural extra basis function at date $k$.
  \item Add $s\mapsto P(t_k,s)$ (notation of \exref{ex:am-stopped-control}) to the four
    Laguerre functions at each date, and redo the experiment: out-of-sample estimate, value of
    the policy, fraction of paths exercised before maturity.
  \item Does the extra function help for the put with $S_0=36$, $\sigma=20\%$, $T=1$ of
    \noteref{tab:am-ls}?
\end{enumerate}
\end{exercise}

\begin{exercise}[Optimising a martingale]\label{ex:am-rogers}
\medium Let $M^E$ be the European put martingale of \noteref{sec:am-duality} and
$H(\lambda)=\E[\max_k(Z_k-\lambda M^E_k)]$ for $\lambda\in\R$.
\begin{enumerate}[label=(\alph*)]
  \item Show that $H$ is convex and that $H(\lambda)\ge U_0$ for every $\lambda$.
  \item For the put with $S_0=36$, $\sigma=20\%$, $T=1$ of \noteref{tab:am-ls}, estimate $H$ on a
    grid of $\lambda\in[0,1.5]$ from $2000$ paths and locate its minimum.
  \item Let $\hat H_n(\lambda)$ be the sample mean from $n$ paths. Show that
    $\E[\min_\lambda\hat H_n(\lambda)]\le\min_\lambda H(\lambda)$. Propose a procedure that
    uses the minimising $\lambda$ and still gives a valid upper confidence bound for $U_0$.
\end{enumerate}
\end{exercise}

\begin{exercise}[The gap of a value-function martingale]\label{ex:am-hk-gap}
\medium Let $\hat V_k$ be functions with $\E[\hat V_k(S_{t_k})^2]<\infty$, set
$\hat U_k=e^{-rt_k}\hat V_k(S_{t_k})$, and define $M_0=0$ and
$M_{k+1}-M_k=\hat U_{k+1}-\E[\hat U_{k+1}\mid\F_k]$, as Haugh and Kogan (2004) do with an
approximate value function. Assume $\E[Z_k^2]<\infty$.
\begin{enumerate}[label=(\alph*)]
  \item Show that $M$ is a square-integrable martingale.
  \item Show that
    $\E[\max_k(Z_k-M_k)]-U_0\le2\big(\sum_{k=1}^m\E[(\hat U_k-U_k)^2]\big)^{1/2}$.
  \item The bound involves the error of $\hat V_k$ everywhere, also far from the exercise
    boundary. Compare with \noteref{prop:am-andersen-broadie}(ii). Which errors of $\hat V_k$
    do not affect $M$ at all?
\end{enumerate}
\end{exercise}

\begin{exercise}[A call on the maximum of two assets]\label{ex:am-max-call}
\hard Let $S^1,S^2$ be independent Black--Scholes stocks with $S^1_0=S^2_0=S_0$,
$\sigma=20\%$ and $y=10\%$, let $r=5\%$, $K=100$, $T=3$, and consider the Bermudan call on
the maximum, with payoff $(\max(S^1,S^2)-K)^+$ and exercise dates $t_k=k/3$, $k=0,\dots,9$.
\begin{enumerate}[label=(\alph*)]
  \item Implement LSM for this option. Choose a basis (for instance polynomials of degree two
    in $S^1/K$ and $S^2/K$, with or without the payoff itself) and justify the choice.
  \item Compute out-of-sample lower bounds for $S_0=90,100,110$ with $10^5$ paths.
  \item Compute Andersen--Broadie upper bounds by nested simulation, with sizes your computer
    allows, and the confidence interval of \noteref{prop:am-andersen-broadie}(iv).
  \item Andersen and Broadie (2004) report the $95\%$ intervals $[8.053,8.082]$,
    $[13.892,13.934]$ and $[21.316,21.359]$ for these three cases. Compute the price itself
    on a two-dimensional grid (the transition factorises over the two independent assets).
    Are your intervals consistent with it? Which half of your interval, below or above the
    price, is wider, and why?
\end{enumerate}
\end{exercise}
